\section*{Project Overview}

Iterative pre-copy live migration transfers memory pages from a running process
to a destination host while the process continues to execute, repeating the
transfer for pages that were dirtied during the previous round until the
remaining working set is small enough to complete migration with a brief final
pause.
The efficiency of this scheme depends entirely on the ability to identify,
at the end of each round, \emph{exactly} which pages were written.
For CPU memory, the Linux kernel provides this capability through soft-dirty
bits embedded in page-table entries, readable via \texttt{/proc/\textit{PID}/pagemap}.
No analogous mechanism exists for GPU-managed memory.
NVIDIA's Unified Virtual Memory (UVM) subsystem, which underpins \texttt{cudaMallocManaged}
and handles demand-paging between host and device, exposes no interface for
querying which pages a process has modified. Consequently, any migration or checkpointing system must conservatively copy the
\emph{entire} GPU memory footprint on every iteration, rendering iterative
pre-copy migration of large GPU workloads impractical.

This project addresses that gap by implementing \textbf{per-process dirty page
tracking} directly within the open-source NVIDIA UVM kernel module.
UVM already intercepts every GPU page fault to manage demand-paging; we augment
the write-fault servicing path to record each written page in an in-kernel
data structure, without introducing any additional hardware overhead.
A \texttt{procfs}-based interface allows a privileged migration daemon to
start and stop tracking epochs and to retrieve the complete set of pages written
by a given process, with \textbf{page-level granularity} and a per-page
\textbf{nanosecond-resolution timestamp} suitable for ordering events across
migration rounds.
